\documentclass[11pt]{article}
\usepackage[a4paper,margin=25mm]{geometry}
\usepackage{amsmath,amssymb,amsthm}
\usepackage[hidelinks]{hyperref}
\emergencystretch=2em
\newtheorem{theorem}{Theorem}
\title{The complement clique number does not bound\newline
the Lov\'asz theta function\\\large Disproof of Conjecture 00000008908}
\author{gaochengzhecpu}\date{}
\begin{document}\maketitle
\begin{abstract}
For the five-cycle $G=C_5$, both its independence number and
the clique number of its complement are two. An explicit
rational Gram certificate for the actual Lov\'asz theta
semidefinite program gives
$\vartheta(G)\geq5041/2257>2$.
Thus these two numbers cannot be the endpoints of the
claimed monotone chain containing $\vartheta(G)$.
Lean checks the full semidefinite feasible set and its
bounded supremum, as well as the actual graph parameters.
\end{abstract}

\section*{The exact endpoint assertion}
The English source explicitly names the endpoints as the
independence number and the \emph{clique number of the
complement}. In particular its proposed sandwich would require
\[
              \alpha(G)\leq\vartheta(G)\leq\omega(\overline G).
\]
The usual sandwich has a different upper endpoint: the
chromatic number $\chi(\overline G)$. The present counterexample
addresses the endpoint actually written in the source.

Let $G$ be the simple cycle on $\mathbb Z/5\mathbb Z$, with
$i$ adjacent to $i+1$ and $i-1$. Its largest independent sets
have size two. Indeed, three independent vertices would
require three cyclic gaps of at least two, impossible on
five vertices; the pair $\{0,2\}$ witnesses size two.
Independent sets of $G$ are precisely cliques of its
complement, so
\[
                   \alpha(G)=\omega(\overline G)=2.
\]
The Lean development proves these parameters using Mathlib's
actual simple graph, complement, clique and independence
definitions, with complete finite certificates.

\section*{The actual theta optimization problem}
Use the standard equivalent semidefinite formulation
\cite[Theorem 4]{lovasz}:
\[
 \vartheta(G)=\sup\left\{\sum_{i,j}X_{ij}:
 X=X^{\mathsf T}\succeq0,\quad\operatorname{tr}X=1,
 \quad X_{ij}=0\text{ if }i\sim_G j\right\}.
\]
The supremum is over \emph{all} feasible real matrices,
rather than only the certificate constructed below.

\newpage
\section*{An exact rational Gram certificate}
Let $B$ be the symmetric circulant matrix with first row
$(-9,32,8,8,32)$, namely
\[
B=\begin{pmatrix}
-9&32&8&8&32\\
32&-9&32&8&8\\
8&32&-9&32&8\\
8&8&32&-9&32\\
32&8&8&32&-9
\end{pmatrix}.
\]
Direct integer multiplication gives
\[
B^{\mathsf T}B=\begin{pmatrix}
2257&0&1392&1392&0\\
0&2257&0&1392&1392\\
1392&0&2257&0&1392\\
1392&1392&0&2257&0\\
0&1392&1392&0&2257
\end{pmatrix}.
\]
Set $X=B^{\mathsf T}B/11285$. It is positive semidefinite
because, for every real vector $z$,
\[
 z^{\mathsf T}Xz=\frac{\lVert Bz\rVert_2^2}{11285}\geq0.
\]
Its trace is $5\cdot2257/11285=1$, and all positions
corresponding to edges of $C_5$ are zero. Thus it is an
actual feasible matrix. Its objective is
\[
 \sum_{i,j}X_{ij}
 =\frac{5\cdot2257+10\cdot1392}{11285}
 =\frac{5041}{2257}
 =2+\frac{527}{2257}>2.
\]

For completeness, the full feasible objective set is bounded.
If $Y\succeq0$ is symmetric, testing on $e_i-e_j$ gives
$2Y_{ij}\leq Y_{ii}+Y_{jj}$. Summing over all $25$ ordered
pairs and using $\operatorname{tr}Y=1$ yields
\[
             2\sum_{i,j}Y_{ij}\leq10,
             \qquad\sum_{i,j}Y_{ij}\leq5.
\]
The certificate makes the feasible set nonempty. Hence
its actual real supremum exists, and it is at least the
certificate's objective. No empty-set or unbounded-set
convention for the supremum is involved.

\newpage
\section*{Contradiction and scope}
\begin{theorem}
For the actual simple five-cycle,
\[
 \alpha(C_5)=\omega(\overline{C_5})=2
       <\frac{5041}{2257}\leq\vartheta(C_5).
\]
Therefore the claimed monotone chain with the two stated
endpoints is false.
\end{theorem}
\begin{proof}
The graph parameters were computed above. The explicit
matrix $X$ is feasible for the full theta optimization
problem and has objective strictly larger than two.
\end{proof}

The proof does not need the familiar exact value
$\vartheta(C_5)=\sqrt5$, and does not claim to formalize
that exact value or all versions of the sandwich theorem.
It proves a sufficient strict rational lower bound against
the exact upper endpoint named in the conjecture.

\section*{Formalization and reproducibility}
The Lean graph is Mathlib's \texttt{cycleGraph 5}.
\texttt{Feasible} uses the actual matrix predicate
\texttt{PosSemidef}, trace one and all actual edge-zero
conditions. \texttt{theta} is the real supremum of the
objective values of every such matrix. The proof establishes
nonemptiness, a bound for the whole feasible set, and the
supremum inequality for the Gram certificate.
The matrix entries and finite graph certificates are exact;
no floating-point optimizer or untrusted numerical oracle
is used.

Lean 4.19.0 and Mathlib revision\\
\texttt{c44e0c8ee63ca166450922a373c7409c5d26b00b}
are pinned. From \texttt{lean/}, run
\begin{verbatim}
lake exe cache get
lake build
lake env lean -DwarningAsError=true Main.lean
\end{verbatim}
The full project is rebuilt afresh and checked with
warnings treated as errors. Printed theorem dependencies
are audited for standard foundational axioms only.
No proof placeholders or custom axioms are used.
Compile the paper with \texttt{tectonic main.tex}.

The exact bilingual source, provenance, actual build logs,
artifact hashes and all-page PDF review are included.
An agent developed the proof and supplied a self-review;
the parent reviewed the complete proof and final artifacts
separately. No external independent review or novelty of
this standard theta counterexample is claimed.

\begin{thebibliography}{9}
\bibitem{source} The Last Math Competition,
Conjecture 00000008908. Exact bilingual source and
provenance included.
\bibitem{lovasz} L. Lov\'asz, \emph{On the Shannon Capacity
of a Graph}, IEEE Transactions on Information Theory,
25(1), 1--7 (1979), Theorem 4.
\href{https://people.orie.cornell.edu/miketodd/orie6327/lovshannon.pdf}{Original-paper mirror}.
This supplies the standard theta definition; the full
certificate needed for this disproof is displayed above.
\end{thebibliography}
\end{document}
